\documentclass[11pt,aps,preprint,nofootinbib,superscriptaddress,tightenlines]{revtex4-2}
\usepackage[T1]{fontenc}
\usepackage{amsmath,amssymb,mathtools}
\let\lambdabar\relax   % revtex defines it; newpxmath does too
\usepackage{newpxtext,newpxmath}
\usepackage{bm}
\usepackage[final]{microtype}
\usepackage{xcolor}
\usepackage{booktabs}
\usepackage{dsfont}
\usepackage[most]{tcolorbox}
\newcommand{\one}{\mathds{1}}
\definecolor{darkblue}{RGB}{20,60,140}
\definecolor{accent}{RGB}{30,80,150}
\definecolor{tint}{RGB}{243,246,251}
\newtcolorbox{keyeq}{enhanced, colback=tint, frame hidden, borderline west={2pt}{0pt}{accent},
  sharp corners, left=4pt, right=4pt, top=0pt, bottom=0pt, before skip=6pt, after skip=6pt}
\usepackage[colorlinks=true,linkcolor=darkblue,citecolor=darkblue,urlcolor=darkblue]{hyperref}

\def\st{\begin{equation}}
\def\stp{\end{equation}}
\def\Eq#1{eq.~(\ref{#1})}
\def\Sect#1{Sect.~\ref{#1}}
\def\Tab#1{Table~\ref{#1}}
\let\oldparagraph\paragraph
\renewcommand{\paragraph}{\par\bigskip\oldparagraph}

\begin{document}
\sloppy

\title{The S-scheme in the code, and its fold at low entropy per particle}
\author{D.~Teaney}
\affiliation{Department of Physics and Astronomy, Stony Brook University,
Stony Brook, New York 11794, USA}
\date{\today}

\begin{abstract}
This note records the work of 2026-11-25 on the state-counting renormalization (the S-scheme, \nolinkurl{renorm3.tex}, \nolinkurl{sscheme_plan.tex}).
The single particle of the S-scheme was compared with split by quadrature; the two maps of the scheme were implemented and tested; the scheme was
put into the steps, which required storing the particle volume.  At $s_a\ge54$ the S-scheme is as good as split or better.  At $s_a=16$ hybrid
Monte Carlo fails: the one-loop Hamiltonian ends on a fold, which the earlier schemes do not have.  We explain why, confirm the explanation by a
test, and describe the regularizations that remove the fold, with the one chosen.  The running record, with every job number, is the Results
section of \nolinkurl{sscheme_plan.tex}.
\end{abstract}

\maketitle
\tableofcontents

%=============================================================================
\section{The scheme}
\label{sec:scheme}

A particle has a label $s$, a momentum $\bm m$ and a volume $V$, with the Liouville measure $e^s\,ds\,d^2m$.  Its \emph{target state} is the temperature
$T_{\rm rel}$ and velocity $\bm v$ that the input equation of state gives to a cell of volume $V$ holding the particle's energy and momentum, and its
\emph{target entropy} is $S_{\rm t}=\gamma\,s(T_{\rm rel})\,V$.  At one loop the S-scheme sets
\begin{keyeq}
\st
s=S_{\rm t}-\Delta,\qquad \Delta=\tfrac D2\ln\big(2\pi V\gamma\,wT_{\rm rel}\big)+\tfrac12\ln a+1+\ln V+\ln\gamma-\ln T_{\rm rel},
\label{scheme}
\stp
\end{keyeq}
with $D=3$ and $a=1/c_s^2-v^2$.  The temperature, the velocity and the pressure of the particle follow by differentiation; at fixed $\bm m$ and $V$,
\st
\mathcal D\equiv\frac{\partial s}{\partial S_{\rm t}},\qquad T_a=\frac{\partial H}{\partial s}=\frac T{\mathcal D} ,
\stp
so a state has a positive temperature only where $\mathcal D>0$.  For a constant speed of sound the counterterm is a sum of three pieces,
\begin{keyeq}
\st
\Delta=b\ln S_{\rm t}+X(\bm m,V)+\text{const},\qquad X=\mathcal N\ln\frac{\gamma V}{\bar V}+\tfrac12\ln\big(1-c_s^2v^2\big),
\label{split}
\stp
\end{keyeq}
with $b=\tfrac D2+1-\mathcal N$, $\mathcal N=1-2c_s^2$.  At $c_s^2=1/2$, $b=5/2$ and $\mathcal N=0$.  The same three pieces are the counterterm of split
(\texttt{SPH\_RENORM=split}): $b\ln S_{\rm t}$ is its label part and $X$ its dynamical part.  The schemes differ only in where they put them.

%=============================================================================
\section{What was done}
\label{sec:done}

\paragraph{The single particle.}  \nolinkurl{examples/renorm_quadrature_sscheme.py} integrates the single-particle partition function of the tree level,
split and the S-scheme, at rest and in a moving bath.  The quadrature had to be moved to the variables $(\ln T,\text{rapidity})$: in $(T,v)$ the weight
of both renormalized models rises toward the light cone, so the Gaussian grid of the earlier script had no peak to sit on and missed a tail worth a few
$10^{-4}$.  Converged to $10^{-6}$, the relative errors of $E$ and $p$ are those of \Tab{tab:single}.  The S-scheme is as good as split or better wherever
one loop applies, by up to a factor of ten in motion; at $s_a=16$ both are off by 4 to 10 percent.

\begin{table}[h]
\caption{The single particle by quadrature (\nolinkurl{examples/renorm_quadrature_sscheme.py}): relative errors of $E$ and of $p$, split and the
S-scheme.  $s_a=16$ is outside the range of one loop for both.}
\label{tab:single}
\small
\begin{tabular}{cccrrrr}
\toprule
 & & & \multicolumn{2}{c}{split} & \multicolumn{2}{c}{S-scheme} \\
$c_s^2$ & $s_a$ & $v$ & $E$ & $p$ & $E$ & $p$ \\
\midrule
$1/2$ & 16 & 0 & $3.9\times10^{-2}$ & $3.9\times10^{-2}$ & $4.2\times10^{-2}$ & $4.2\times10^{-2}$ \\
 & 54 & 0 & $2.4\times10^{-4}$ & $2.4\times10^{-4}$ & $1.8\times10^{-4}$ & $1.8\times10^{-4}$ \\
 & 54 & 0.8 & $6.9\times10^{-4}$ & $8.0\times10^{-4}$ & $2.8\times10^{-4}$ & $2.4\times10^{-4}$ \\
 & 541 & 0.8 & $6.1\times10^{-6}$ & $7.2\times10^{-6}$ & $2.3\times10^{-6}$ & $2.0\times10^{-6}$ \\
$1/4$ & 54 & 0 & $5.3\times10^{-4}$ & $2.4\times10^{-3}$ & $5.8\times10^{-4}$ & $5.7\times10^{-4}$ \\
 & 54 & 0.8 & $-4.4\times10^{-4}$ & $1.0\times10^{-3}$ & $4.4\times10^{-5}$ & $-1.0\times10^{-4}$ \\
 & 541 & 0.8 & $-3.7\times10^{-6}$ & $9.5\times10^{-6}$ & $-1.3\times10^{-7}$ & $-1.5\times10^{-6}$ \\
\bottomrule
\end{tabular}
\end{table}

\paragraph{The maps.}  \nolinkurl{src/sscheme.jl} evaluates everything about a particle at its target state: the label, the energy, $T_a$, $\bm v_a$, $P_a$,
$\mathcal D$ and the weight of a vote, $s-\ln T_a$, which needs no solve.  For a constant speed of sound both target inversions are one-dimensional.
The map from the label iterates on $S_{\rm t}$ from its previous value (1.7 Newton steps per solve in hybrid Monte Carlo).  The tests
(\nolinkurl{test/testsscheme.jl}) check $\partial H/\partial s=T_a$, $\partial H/\partial\bm m=\bm v_a$, $-\partial H/\partial V=P_a$ to $3\times10^{-10}$, the
two maps against each other to $2\times10^{-15}$, an independent implementation in mpmath to $10^{-15}$, and the relation \Eq{split} to split.

\paragraph{The steps: the volume is stored.}  The code kept the volume as the smoothed entropy density $S_a=s_an_a$ and recovered $V_a=s_a/S_a$, the
energy $s_a\,(H/s)_a$ and the velocity $(v/S)_aS_a$ through the label.  The label of the S-scheme can be zero or negative, so the particle now stores
its volume, energy, velocity and target entropy.  Every scheme fills them, the tree level, split and Legendre with the expressions used before, so
their dynamics is bit-identical to \texttt{main}.  Three diagnostics of the equilibrium driver changed because they now read a consistent state: the
energy series (the old expression mixed labels after the stochastic step with energies before it, a noise of $10^{-5}$ in $\Delta H/H$; now $10^{-7}$),
and the kinetic energy and the volume spread.  With the S-scheme the ideal step and hybrid Monte Carlo use the particle's own $T_a$, $\bm v_a$, $P_a$;
the votes use $s-\ln T_a$; the transport coefficients and the comoving structures of the viscous step use the target state.  Hybrid Monte Carlo
conserves $H$ to second order in $\delta t$ (ratio 16.1 under two halvings), and the full test suite passes.

\paragraph{The dynamics at $s_a=54$.}  With the S-scheme, at rest, 5 time units: hybrid Monte Carlo acceptance 0.99, $\Delta H/H=-5.9\times10^{-6}$,
KE/eq $0.984$; split $-2.2\times10^{-6}$, $0.986$.  At $v_0=0.8$, $\Delta H/H=+3.5\times10^{-7}$.  The equilibrium campaign at $s_a=54$ and $541$
(\texttt{eos\_sscheme}) is running.

%=============================================================================
\section{The fold at $s_a=16$}
\label{sec:fold}

\paragraph{What was seen.}  At $s_a=16$, $c_s^2=1/2$, at rest, hybrid Monte Carlo accepted 11 percent of its trajectories (split 99 percent), and the
dynamics drifted by $1.9\times10^{-3}$ in 5 time units (split $3\times10^{-4}$).  Of 40 trajectories, 32 left the domain of $H$ and none was
rejected by Metropolis; at a time step 3.5 times smaller, 31 left.  The trajectories that stayed conserved $H$ to $10^{-6}$.  The exits are a property
of $H$, not of the integrator.

\paragraph{Why.}  At fixed $\bm m$ and $V$, \Eq{split} gives
\st
s=S_{\rm t}-b\ln S_{\rm t}-X(\bm m,V)-\text{const},\qquad \frac{\partial s}{\partial S_{\rm t}}=1-\frac b{S_{\rm t}} .
\stp
The label has a minimum at $S_{\rm t}=b$, the fold, where $\mathcal D=0$ and $T_a$, $v_a$ diverge.  Its value,
\st
s_{\min}(\bm m,V)=b-b\ln b-X(\bm m,V)-\text{const},
\stp
depends on $\bm m$ and $V$ through $X$.  The ideal step conserves the label while $\bm m$ and $V$ change.  If $X$ falls, $s_{\min}$ rises, and a particle
whose label lies just above $s_{\min}$ is left with no state: it has run onto the fold.  At $c_s^2=1/2$, $X=\tfrac12\ln(1-v^2/2)$ can fall by $\tfrac12\ln2=0.35$,
so a particle within $0.35$ of the minimum can be carried onto the fold by its own acceleration.  At $s_a=16$ about one percent of the particles are
that close at any moment; at $s_a=54$ essentially none are.

\paragraph{Why split and \texttt{SPH\_RENORM=1} do not have it.}  They take the target entropy $u=S_{\rm t}$ as the particle's variable.  The label part
$b\ln u$ sits in the measure, $e^{u-b\ln u}(1-b/u)\,du$, and their boundary $u>b$ is on the variable that the ideal step conserves, so it cannot be
crossed.  The part $X$ is left out (\texttt{SPH\_RENORM=1}) or added to $H$ as an energy (split), where it moves nothing.  The S-scheme puts $X$ inside
the conserved label.

\paragraph{The test.}  With the sound term removed from $\Delta$, the S-scheme at $c_s^2=1/2$ has $X=0$, and its ideal step is that of
\texttt{SPH\_RENORM=1}.  Then 37 of 40 trajectories were accepted and none left the domain.

\paragraph{It belongs to the truncation.}  The exact number of states increases with the energy at every $(\bm m,V)$, and at $\bm m\ne0$ its labels cover
the real line (\nolinkurl{renorm2.tex}): the exact scheme has no fold.  The next term of the series does not remove it.  With $\Delta\to\Delta-c/S_{\rm t}$ the
slope $1-b/S_{\rm t}-c/S_{\rm t}^2$ still vanishes, near $S_{\rm t}\approx b+c/b$, and $c$ depends on $v$.  The fold sits where the expansion parameter is of
order one: the exact particle is a Poisson cloud of $\bar N=S_{\rm t}/n$ quanta, and at the fold $\bar N\approx1$.

%=============================================================================
\section{Removing the fold}
\label{sec:fix}

Two conditions are needed.  At fixed $\bm m$ and $V$ the label must increase with $S_{\rm t}$, and its smallest value must not depend on $\bm m$ and $V$.
One loop has the first and not the second.  Any change confined to $S_{\rm t}\sim b$, where one loop does not apply, is admissible, and changes of order
$1/S_{\rm t}^2$ are accepted.

\paragraph{A hard limit on the label.}  Allow only $s>s_b=\max_{\bm m,V}s_{\min}$; for $c_s^2=1/2$, $s_b=s_{\min}(\text{rest})+\tfrac12\ln2$.  The ideal step
conserves $s$ and cannot reach the fold.  The stochastic step, the heat bath of hybrid Monte Carlo and the equilibrium draw are the only places where
labels change, and they reject a state below the limit as they reject any state outside the domain.  It removes the labels within $0.35$ of the minimum
(at rest the target entropies $b<S_{\rm t}\lesssim4$) and leaves $H$ unchanged elsewhere.  The volume is not touched.

\paragraph{Switching off $X$ at the boundary.}  Replace $X$ by $g(S_{\rm t})X$ with
\st
g=\Big(1-\frac b{S_{\rm t}}\Big)^2\Big(1+\frac{2b}{S_{\rm t}}\Big)=1-\frac{3b^2}{S_{\rm t}^2}+\frac{2b^3}{S_{\rm t}^3} .
\stp
Since $g(b)=g'(b)=0$ the smallest label is the same for every $\bm m$ and $V$ (checked: $-2.79548$ at $\gamma v=0,1,3,10$), and the change is of order
$1/S_{\rm t}^2$.  The slope becomes $(1-b/S_{\rm t})(1-6b^2X/S_{\rm t}^3)$, so the second factor must stay positive, $X<S_{\rm t}^3/6b^2\approx0.33$ near the
boundary.  It also removes a symptom of the fold: a particle at $S_{\rm t}=3$ with $|\bm m|=2$ moved against its momentum ($v_a=-0.15$) and now moves with it
($v_a=0.57$).

\paragraph{The choice (owner).}  Multiply by $g$ only the parts of $X$ that raise the smallest label, and put the hard limit at the rest minimum, as split's
$u>b$.  The velocity part, $\mathcal N\ln\gamma+\tfrac12\ln(1-c_s^2v^2)$, raises the minimum for $c_s^2>1/3$ and lowers it for $c_s^2<1/3$.  The volume part,
$\mathcal N\ln(V/\bar V)$, raises it for $V<\bar V$ whenever $\mathcal N>0$.  So $g$ multiplies the velocity part at $c_s^2=1/2$ and the volume part at
$c_s^2=1/4$.  Putting $g$ on every term would need $g'X<1-b/S_{\rm t}$, which the velocity part at $c_s^2=1/4$ violates for fast particles near the
boundary ($\tfrac12\ln\gamma$ reaches 2.3 at $\gamma=100$).  The volume part with $g$ needs $V<1.9\bar V$ near the boundary, five standard deviations in the
measured spread of the volumes at $s_a=16$.

\paragraph{The sign of the particle pressure (owner's question, 2026-11-25).}  The fold is the statement $T_a=T/\mathcal D>0$, with $T_a\to\infty$ at the
fold and $T_a<0$ beyond it.  The pressure is a separate condition.  At one loop, for a constant speed of sound,
\st
\frac{P_aV}{T_a}=\frac{S_{\rm t}}n+\frac{\tfrac D2(2+q)-1}{n}-\Big(\frac D2+1\Big),
\stp
the same at every $v$ and $V$, and the constant is negative: $-0.83$ at $c_s^2=1/2$, $-0.90$ at $c_s^2=1/4$.  At $c_s^2=1/2$ the pressure is positive
exactly above the boundary, $S_{\rm t}>5/2=b$.  At $c_s^2=1/4$ it is negative for $b=2<S_{\rm t}<4.5$: particles just above the boundary pull their
neighbours in.  With the regularization the factor $g$ on the volume part adds $\partial(hX_{\rm vol})/\partial\ln V$ terms; at $V=\bar V$ the
addition is $h\mathcal N$, $+\tfrac12$ at the boundary, which cancels the deficit there.  $P_aV/T_a$ on a grid ($A=1$ units;
\texttt{SPH\_SSCHEME\_REG}$=0/1$; -- outside the domain):
\begin{center}\small
\begin{tabular}{lcrrrrrr}
\toprule
 & & \multicolumn{6}{c}{$S_{\rm t}$} \\
 & $v$, $V/\bar V$ & $b{+}0.1$ & $b{+}0.5$ & 3.5 & 4.5 & 6 & 10 \\
\midrule
$c_s^2=1/2$, one loop and regularized & $0$, any & $+0.033$ & $+0.167$ & $+0.333$ & $+0.667$ & $+1.167$ & $+2.500$ \\
\quad regularized & $0.8$, any & $+0.047$ & $+0.211$ & $+0.390$ & $+0.720$ & $+1.206$ & $+2.518$ \\
$c_s^2=1/4$, one loop & any & $-0.480$ & $-0.400$ & $-0.200$ & $0.000$ & $+0.300$ & $+1.100$ \\
\quad regularized & $0$, $0.5$ & $+0.035$ & $+0.101$ & $+0.161$ & $+0.254$ & $+0.460$ & $+1.165$ \\
 & $0$, $1$ & $+0.017$ & $+0.048$ & $+0.103$ & $+0.209$ & $+0.430$ & $+1.152$ \\
 & $0$, $1.5$ & $+0.006$ & $+0.017$ & $+0.069$ & $+0.182$ & $+0.412$ & $+1.144$ \\
 & $0.8$, $1$ & -- & -- & $+0.103$ & $+0.209$ & $+0.430$ & $+1.152$ \\
\bottomrule
\end{tabular}
\end{center}
So the regularization chosen also makes the particle pressure positive above the boundary on this grid, at both speeds of sound.  At large $V$ the
margin shrinks ($+0.006$ at $S_{\rm t}=b+0.1$, $V=1.5\bar V$), because the term $\partial h/\partial\ln S_{\rm t}\,X_{\rm vol}/n$ is negative for $V>\bar V$;
whether it turns negative for still larger volumes is to be checked with the cutoff family (E2).

\paragraph{The tests at $s_a=16$, $c_s^2=1/2$} (40 updates of hybrid Monte Carlo, and 5 time units of the dynamics):
\begin{center}\small
\begin{tabular}{lcccc}
\toprule
 & accepted & left the domain & $\Delta H/H$ & KE/eq \\
\midrule
one loop & 8 & 32 & $+1.9\times10^{-3}$ & 1.003 \\
hard limit & 40 & 0 & & \\
$g$ & 39 & 0 & $-9.8\times10^{-6}$ & 0.996 \\
split & 40 & 0 & $-3.4\times10^{-6}$ & 1.024 \\
\bottomrule
\end{tabular}
\end{center}
The single particle is not affected where one loop applies.  With the hard limit the errors at $s_a=54$ are identical to one loop's.  With $g$ they are
unchanged at $s_a=541$; at $s_a=54$, $v=0.8$ they grow from $2.8\times10^{-4}$ to $4.2\times10^{-4}$, because $g$ approaches one as $1/S_{\rm t}^2$.  At $s_a=16$
all three are near 4 percent, as split is.  The guard rejections of the stochastic step at $s_a=16$ (about $1.2\times10^5$ in 5 time units) are the same
for $g$ and split, so they are a property of the entropy per particle, not of the fold.

%=============================================================================
\section{Where things stand}
\label{sec:next}

\begin{itemize}
\item Merged into the integration branch \texttt{sscheme}: the maps, the steps, the tests.  The regularizations are diagnostic settings on the branch
  \texttt{sscheme-fold-diag}, for $c_s^2=1/2$ only.
\item Running: the equilibrium campaign \texttt{eos\_sscheme} at $s_a=54$ and $541$ (\nolinkurl{sscheme_plan.tex}, item F).
\item Next: the dynamical campaign at $s_a=54$ (item G); the chosen regularization implemented as a setting with its tests, and checked at $c_s^2=1/4$;
  then the family of cutoff functions and the sensitivity of the results at $s_a\lesssim30$ to them (item E2).
\item Open, for the owner: a label written as a function of $\Delta/S_{\rm t}$ that is increasing for every $S_{\rm t}$, guided by the known higher orders and the
  exact Poisson cloud.
\end{itemize}

\end{document}
